\section{Conclusion}
\label{sec:conclusion}

When accounts cannot be linked, what a provider can learn from queries alone is decided by benign coordination: a campaign is visible only if it stands out more than the most distinctive group of ordinary users. That one idea explains why naive distributed extraction is catchable, why padding blinds detectors that weigh accounts by size but not one that weighs them by direction, why selection fails before the barrier as populations grow, and why relays are the hard case.

For a provider, the result is a filter worth running and a clear list of what it will not do. A calibrated, hollowed monitor catches dataset-like campaigns spread across many accounts, keeps catching them when they are padded, and points reviewers at a short list of accounts that are mostly colluders; in our testbed, the campaigns it could not see had given up most of their value. But it detects topic shift, not extraction; a slow ramp badly degrades it; we found no deployable configuration at production colluder shares; and relays may be beyond any passive monitor. The decisive quantity---how coordinated a provider's own benign traffic is---can only be measured on real logs, and we hope the census script, the calibration protocol and the record of what failed make that measurement easy to repeat.
